\section{EXPERIMENTS}

\subsection{Experimental Setup}
\label{sec:setup}

\paragraph{Simulation data.}
We generate paired RGB-tactile training data in Blender for 20 objects across
12 sessions per object. Each sample consists of a tactile gel image, an
external RGB image, ground-truth depth (float32 \texttt{.npy}), and a
6-DOF pose annotation. A fixed $1/\sqrt{2}$ center crop yields an effective
$12.37$\,mm tactile field of view. All images are resized to
$224{\times}224$.

\paragraph{Real data.}
\todo{Describe your real sensor platform, data collection (5 objects,
session\_000), and how real samples are formatted to match sim.}

\paragraph{Metrics.}
We report: depth MAE and RMSE (contact region), normal angular error
(mean, degrees), pose rotation error (degrees, via
$\arccos(\cos\hat\theta\cos\theta^* + \sin\hat\theta\sin\theta^*)$),
and pose translation L1. All metrics are computed on the simulation
validation set (2400 samples, every 20th sample per session) unless
otherwise noted.

\paragraph{Baselines.}
We compare against: (1)~a shared-encoder baseline using a single frozen
DINOv3~\cite{simeoni2025dinov3} for both modalities with symmetric MBT-style
fusion; (2)~MViTac~\cite{fu2024mvitac}, a contrastive ResNet-18 baseline;
(3)~VITaL, a CLIP-pretrained ResNet-18 baseline with DPT decoder.
All baselines are trained with the same data, augmentation, and loss
configuration.

% -------------------------------------------------------------------------
\subsection{Does Fusion Beat Either Modality Alone?}
\label{sec:modality_ablation}

The central hypothesis is tested with the fair single-model protocol:
identical trainable parameters and identical decoder input across all three
configurations.

\begin{table}[t]
\caption{Modality ablation (one shared model, T3+MAE encoders, sim-only
training). Dense prediction is driven by tactile; pose by RGB; fusion
recovers both.}
\label{tab:modality_ablation}
\begin{center}
\begin{tabular}{lcccc}
\toprule
Config & \begin{tabular}{@{}c@{}}Depth\\MAE$\downarrow$\end{tabular}
       & \begin{tabular}{@{}c@{}}Normal\\(deg)$\downarrow$\end{tabular}
       & \begin{tabular}{@{}c@{}}Rot.\\(deg)$\downarrow$\end{tabular}
       & \begin{tabular}{@{}c@{}}Trans.\\L1$\downarrow$\end{tabular} \\
\midrule
RGB-only      & 0.496  & 12.67 & 2.18  & 0.008 \\
Tactile-only  & 0.013  & 1.21  & 89.94 & 0.205 \\
RGB+tactile   & \textbf{0.013} & \textbf{1.18}  & \textbf{2.21}  & \textbf{0.008} \\
\bottomrule
\end{tabular}
\end{center}
\end{table}

Table~\ref{tab:modality_ablation} reveals a striking complementarity.
Tactile-only achieves near-identical depth and normal accuracy as the fused
model (MAE 0.013, $1.21^{\circ}$), confirming that contact geometry is
overwhelmingly a tactile signal. However, tactile-only pose estimation is at
chance ($89.94^{\circ}$ $\approx$ uniform over $[-180^{\circ},
180^{\circ}]$), because geometrically symmetric objects cannot be
disambiguated from gel deformation alone.
Conversely, RGB-only recovers pose ($2.18^{\circ}$) via visual object
identity but cannot reconstruct contact depth (MAE $38{\times}$ worse).
Fusion inherits the strengths of both: the dense accuracy of tactile and
the pose disambiguation of RGB.

% -------------------------------------------------------------------------
\subsection{Encoder Ablation}
\label{sec:encoder_ablation}

We compare six frozen encoder pairings (tactile $\times$ RGB) while keeping
all other hyperparameters fixed. Training includes sim+real co-training.

\begin{table}[t]
\caption{Encoder ablation (sim+real co-training). Rotation and depth
evaluated in ``both'' (RGB+tactile) mode.}
\label{tab:encoder_ablation}
\begin{center}
\begin{tabular}{llcc}
\toprule
Tactile enc. & RGB enc. & Rot.\ (deg)$\downarrow$ & Depth MSE$\downarrow$ \\
\midrule
\textbf{T3}     & \textbf{MAE}    & \textbf{0.803} & 0.0195 \\
DINOv3 (shared) & DINOv3 (shared) & 0.812          & 0.0204 \\
SITR            & MAE             & 0.819          & \textbf{0.0159} \\
TVL             & MAE             & 0.844          & 0.0223 \\
T3              & DINOv3          & 1.481          & 0.0204 \\
T3              & CLIP            & 1.659          & 0.0208 \\
\bottomrule
\end{tabular}
\end{center}
\end{table}

Table~\ref{tab:encoder_ablation} shows that T3+MAE achieves the best pose
accuracy ($0.803^{\circ}$), outperforming the single shared DINOv3 baseline
($0.812^{\circ}$) despite using two separately pretrained encoders.
SITR+MAE achieves the best depth MSE, suggesting that calibration-aware
tactile encoders may better capture absolute surface geometry. The visual
encoder choice matters substantially for pose: MAE consistently outperforms
CLIP and DINOv3 as the RGB backbone, likely because its patch-level
reconstruction pretext captures the spatial appearance details needed for
orientation estimation.

% -------------------------------------------------------------------------
\subsection{Comparison with Baselines}
\label{sec:baselines}

\begin{table}[t]
\caption{Comparison with baselines (sim+real co-training, ``both'' mode).}
\label{tab:baselines}
\begin{center}
\begin{tabular}{lcc}
\toprule
Model & Rot.\ (deg)$\downarrow$ & Depth MSE$\downarrow$ \\
\midrule
\textbf{\method{} (T3+MAE)} & \textbf{0.803} & \textbf{0.0195} \\
Shared DINOv3 (MBT)         & 0.812          & 0.0204 \\
MViTac~\cite{fu2024mvitac}  & 1.203          & 0.1529 \\
VITaL                       & 2.271          & 0.0270 \\
\bottomrule
\end{tabular}
\end{center}
\end{table}

Table~\ref{tab:baselines} compares \method{} against three baselines trained
with the same data and losses. MViTac, which uses a contrastive ResNet-18,
achieves reasonable pose but poor depth reconstruction (MSE $7.8{\times}$
higher), highlighting the advantage of ViT-based multiscale features for
dense prediction. VITaL improves depth via its DPT decoder but trails in pose
accuracy ($2.271^{\circ}$ vs.\ $0.803^{\circ}$). The shared-DINOv3 baseline
with symmetric MBT fusion is a strong competitor; \method{}'s advantage comes
from the asymmetric design and modality-matched encoders.

% -------------------------------------------------------------------------
\subsection{Sim-to-Real Transfer}
\label{sec:sim2real}

\paragraph{Role of real data.}
Table~\ref{tab:real_ratio} reports the effect of varying the amount of real
data during co-training (with a fixed sim dataset).

\begin{table}[t]
\caption{Real data quantity ablation (T3+MAE, ``both'' mode).}
\label{tab:real_ratio}
\begin{center}
\begin{tabular}{rcc}
\toprule
Real samples & Rot.\ (deg)$\downarrow$ & Depth MSE$\downarrow$ \\
\midrule
0 (sim-only)    & 3.83  & \todo{--}  \\
25              & 2.126 & 0.0622 \\
100             & 1.029 & 0.0408 \\
300             & 0.825 & 0.0236 \\
all             & 0.886 & 0.0220 \\
\bottomrule
\end{tabular}
\end{center}
\end{table}

Performance improves rapidly with the first hundred real samples and saturates
around 300, with diminishing returns (and slight overfitting) beyond. The
optimal real-to-sim ratio is approximately 1:2.7.
\todo{Expand with your hardware/data-collection perspective on what makes
this ratio work (diversity of real contacts, etc.).}

\paragraph{Domain adaptation.}
We also evaluated explicit domain adaptation methods---CORAL feature alignment
and Fourier Domain Adaptation (FDA)---but neither improved over the
frozen-encoder baseline ($0.846^{\circ}$ and $0.845^{\circ}$ vs.\
$0.803^{\circ}$). This suggests that the frozen pretrained encoders already
provide sufficient domain invariance, and sim+real co-training handles the
residual gap.

\paragraph{Real sensor evaluation.}
\todo{Qualitative and quantitative results on real sensor data. Show
depth/normal visualizations and pose accuracy on real objects.}

% -------------------------------------------------------------------------
\subsection{Qualitative Results}
\label{sec:qualitative}

\todo{Figure: grid of depth/normal/pose predictions for representative
objects across the three input configurations (both, tactile-only, rgb-only),
on both sim and real data. Include failure cases (e.g.\ symmetric objects
under tactile-only).}
